\documentclass[10pt,a4paper]{amsart}
\usepackage[margin=19mm]{geometry}
\usepackage{amsmath,amssymb,mathtools}
\usepackage[scr=rsfs,cal=euler]{mathalfa}
\usepackage{xfrac}
\usepackage[unicode,hidelinks,bookmarks=false]{hyperref}
\newcommand{\ie}{i.e.\ }
\newcommand{\cf}{cf.\ }
\newcommand{\resp}{respectively\ }
\newcommand{\Subsection}{\subsection}
\providecommand{\nobreakdash}{\nobreak\hbox{-}}
\setlength{\emergencystretch}{2em}
\multlinegap0pt
\usepackage{bm}
\usepackage{bbm}
\usepackage{dsfont}
\usepackage{xspace}
\usepackage{cancel}
\usepackage{mathtools}
\usepackage{cleveref}
\usepackage{amsthm}
\makeatletter
\@ifundefined{theorem}{\newtheorem{theorem}{Theorem}}{}
\@ifundefined{lemma}{\newtheorem{lemma}[theorem]{Lemma}}{}
\makeatother
\newcommand{\Rcop}{R_{\mathrm{cop}}}
\title[Prime Certificates for Exact Vertex-Coprime Ramsey Numbers]{Prime Certificates for Exact Vertex-Coprime Ramsey Numbers}
\author{Zhicheng Du, Wenji Xi, Zhuo Deng, Lan Ma}
\date{Source: arXiv:2605.26815v2 (CC BY 4.0)}
\begin{document}
\maketitle

\noindent\emph{Source and licence.} Prime Certificates for Exact Vertex-Coprime Ramsey Numbers. Version arXiv:2605.26815v2 (CC BY 4.0), distributed under the Creative Commons Attribution 4.0 International licence, as recorded in the arXiv licence metadata for this exact version and shown on the abstract page https://arxiv.org/abs/2605.26815v2. Full rights notice: licenses/arxiv-2605.26815-CC-BY-4.0.txt.\par\smallskip

\noindent\textit{Editorial scope.} Counted unit: the Theorem whose source label is \texttt{thm:balanced-exact} in the article (source0.tex lines 1027--1035, source label \texttt{thm:balanced-exact}), delivered together with the article's own complete original proof of it (source0.tex lines 1037--1090, one \texttt{proof} environment of 54 lines). No mathematics is written, repaired or re-derived: statement and proof are the article's own text line for line. Required context retained uncounted: thm:main (lines 226--234); lem:prime-index-gaps (lines 996--1002). Every definition, display and label of those lines is retained unchanged. Omitted from this package and not redistributed: 1--225, 235--995, 1003--1026, 1036--1036, 1091--2900. Nothing inside a retained excerpt is omitted or altered: every retained body is delivered exactly as the article's lines are written, so no omission receipt of any kind accompanies this package. Citations: the retained excerpts carry no citation command at all, so no bibliography entry of the article is transcribed and no reference is redistributed. Reference adaptation: none was needed and none was made; every reference inside the delivered unit resolves inside it, so no unresolved reference or citation is produced. Printed slips kept literal: no printed word, symbol, hypothesis or number of the counted statement or of its proof was altered. Layout and class: the article's own class is \texttt{article} with packages including geometry, amsmath, amssymb, amsthm, array, booktabs, graphicx, xurl, hyperref, cleveref, microtype; the wrapper is the lane's pinned \texttt{amsart} editorial wrapper at 10pt on A4, which restates the theorem-like environments the retained text needs and adds the article's own macro definitions that the retained text needs (\texttt{\textbackslash Rcop}), with extra wrapper packages (bm, bbm, dsfont, xspace, cancel, mathtools, cleveref). The delivered numbering is the unit's own. Assets: no retained span contains a figure environment, a graphics-inclusion command, a PDF-page inclusion, a table or a TikZ picture; no image file is copied into this package, the package declares no asset dependency and no image file is redistributed with it. Licence: version arXiv:2605.26815v2 of this article is distributed under the Creative Commons Attribution 4.0 International licence, as recorded in the arXiv licence metadata for this exact version and shown on the abstract page https://arxiv.org/abs/2605.26815v2. A full rights notice is delivered at licenses/arxiv-2605.26815-CC-BY-4.0.txt. Characters: every retained excerpt is pure ASCII, so the delivered engine, XeLaTeX with its default Latin Modern text font, reports no missing character.
\begin{theorem}[Exact mixed formula]
\label{thm:main}
For all $c\ge 1$ and $k_1,\ldots,k_c\ge 2$,
\[
  \Rcop(k_1,\ldots,k_c)=p_M,
  \qquad
  M=\sum_{i=1}^c(k_i-1).
\]
\end{theorem}

\begin{lemma}[Weak prime-index gaps]
\label{lem:prime-index-gaps}
For every $m\ge2$,
\[
  2p_m<p_{2m}<3p_m .
\]
\end{lemma}

\begin{theorem}[Exact balanced endpoint]
\label{thm:balanced-exact}
For every $k\ge2$,
\[
  L_{\mathrm{bal}}(k;2)=p_{2k-2}-1.
\]
Equivalently, the balanced upper transition point is $p_{2k-2}$, the same as
the unrestricted two-color threshold.
\end{theorem}

\begin{proof}
The upper bound follows from \Cref{thm:main}.  The case $k=2$ is immediate
(place $1$ and $2$ in opposite colors), so assume $k\ge3$ and put
$x=p_{2k-2}-1$.  Then $x$ is even and the primes at most $x$ are exactly
$p_1,\ldots,p_{2k-3}$.

\medskip\noindent\textit{The deterministic split.}
Let
\[
  B_0=\{p_2,\ldots,p_{k-1}\},\qquad
  B_1=\{p_1,p_k,\ldots,p_{2k-3}\},
\]
so $|B_0|=k-2$ and $|B_1|=k-1$.  We will produce a two-coloring with vertex
$1$ in color $0$ and with every other vertex $v$ assigned color $i$ having
at least one prime divisor in $B_i$.  Such a coloring is divisor-certified:
color $0$ has clique size at most $1+|B_0|=k-1$ (vertex $1$ plus injected
witness primes) and color $1$ has clique size at most $|B_1|=k-1$, so no
monochromatic coprime $K_k$ exists.

\medskip\noindent\textit{Vertices forced into color $0$.}
Write $\sigma(v)$ for the set of prime divisors of $v$, and define
\[
  F_0=\{1\}\cup\{v\in[2,x]:\sigma(v)\subseteq B_0\}.
\]
Vertices in $F_0$ must be colored $0$ to preserve the certificate.  Since
$B_0$ contains only odd primes, every $v\in F_0$ is odd.  Conversely, an odd
$v\in[1,x]$ lies in $F_0$ unless $v$ has a prime divisor $q\ge p_k$; if
additionally $v\ne q$, then $v$ has a second prime factor of size at least
$3$, so $v\ge 3q\ge 3p_k>p_{2k-2}=x+1$ by \Cref{lem:prime-index-gaps}, a
contradiction.  The odd vertices in $[1,x]$ outside $F_0$ are therefore
exactly the primes $p_k,\ldots,p_{2k-3}$, of which there are $k-2$.  Since
$[1,x]$ contains $x/2$ odd integers,
\[
  |F_0|=\frac{x}{2}-(k-2).
\]

\medskip\noindent\textit{Flexible vertices moved into color $0$.}
For $i=2,\ldots,k-1$, set $f_i=2p_i$.  By \Cref{lem:prime-index-gaps},
$2p_{k-1}<p_{2(k-1)}=p_{2k-2}=x+1$, so every $f_i\le x$.  The support
$\sigma(f_i)=\{2,p_i\}$ meets both bins: $2\in B_1$ and $p_i\in B_0$, so
$f_i$ may be colored either way without breaking the certificate.  Assign
these $k-2$ vertices to color $0$.

\medskip\noindent\textit{Assigning the remaining vertices.}
Color every remaining vertex with $1$.  Each such vertex either has support
in $B_1$ (forced into color $1$) or has support meeting $B_1$ (flexible,
not selected above); in both cases coloring it $1$ uses a prime in $B_1$,
so the certificate is preserved.

\medskip\noindent\textit{Balance.}
Color $0$ now has $|F_0|+(k-2)=x/2$ vertices and color $1$ has the remaining
$x/2$.  The two-coloring is therefore exactly balanced and avoids
monochromatic coprime $K_k$ in both colors.
\end{proof}

\end{document}
